\section{Síntesis y ampliación}

Las dos conferencias convergen finalmente en una misma disciplina de investigación: no mistificar las capacidades, las arquitecturas ni las tendencias; primero leer los datos y las gráficas, descomponer las tareas y los supuestos ocultos, y después buscar los puntos de cruce a lo largo del eje de escala. Jason ofrece una descomposición microscópica del comportamiento del modelo y Hyung Won, una descomposición macroscópica de la evolución de las arquitecturas; ambos exigen al investigador distinguir entre observación, explicación y extrapolación.

\subsection{Doce conclusiones centrales}

Esta sección condensa las extensas notas en una lista de verificación reutilizable. Cada conclusión conserva sus condiciones de aplicación, para evitar convertir un juicio de clase de 2024 en un hecho de 2026 sin límites.

\begin{enumerate}
  \item Revisar los datos a mano entrena la representación que el propio investigador tiene de la tarea y no puede sustituirse por completo con un benchmark dashboard.
  \item La next-token prediction reúne, mediante continuations reales, una enorme cantidad de latent-task supervision.
  \item La cross-entropy global es una mezcla ponderada de distintas distribuciones de tokens y dificultades de tarea.
  \item Que la loss global sea suave no exige que cada benchmark metric lo sea.
  \item Una emergent-looking curve puede contener a la vez un cambio real de capacidad, una puntuación por umbral y un muestreo disperso.
  \item El inverse/U-shaped scaling suele indicar estrategias en competencia; conviene descomponerlo en subtareas y medirlas.
  \item Una mejora en un solo punto no muestra si el método se satura, si desplaza la curva o si cambia su pendiente.
  \item El cómputo efectivo cada vez más barato es una tendencia histórica sólida, pero no una garantía ilimitada.
  \item El inductive bias mejora la eficiencia en el regime actual, pero puede convertirse en un cuello de botella a mayor escala.
  \item La diferencia entre encoder-decoder y decoder-only puede descomponerse en cuatro hipótesis locales que admiten ablation.
  \item La translation y las tareas de entrada larga con objetivo corto explican por qué ciertas estructuras fueron eficaces en su momento.
  \item La próxima estructura que habrá que reexaminar puede estar en el objective, en la interfaz de datos o en la interacción del sistema, y no solo en el architecture block.
\end{enumerate}

\subsection{Una tabla que conecta a los dos ponentes}

Para situar el comportamiento microscópico y la arquitectura macroscópica en un mismo marco, la siguiente tabla alinea las dos conferencias según «objeto de observación, unidad de descomposición, eje de escala, riesgo y acción». La tabla no presenta conclusiones nuevas; ayuda al lector a trasladar el método.

\begin{table}[H]
\centering
\begin{tabularx}{\textwidth}{L{0.18\textwidth} X X}
\toprule
Dimensión & Jason Wei & Hyung Won Chung \\
\midrule
Objeto de observación & token, prompt, task metric, loss curve & arquitectura, tareas históricas, compute trend, system cost \\
Unidad de descomposición & latent task, subskill, metric threshold & parameter sharing, layer connection, attention mask, objective \\
Eje de escala & model/data/compute budget & compute availability, network depth, interaction length \\
Riesgo principal & Tomar la puntuación agregada por un mecanismo y un salto abrupto por una capacidad misteriosa & Volver permanentes los atajos históricos y tomar los benchmarks antiguos por las tareas futuras \\
Acción recomendada & Mirar los datos, descomponer las tareas, trazar scaling curves con varios puntos & Identificar la fuerza dominante, registrar los sesgos, hacer ablations por sustracción \\
\bottomrule
\end{tabularx}
\caption{Correspondencia entre las dos conferencias dentro de un mismo marco guiado por la evidencia.}
\end{table}

\subsection{Plantilla de investigación para la práctica}

El lector puede aplicar el método del curso directamente a una idea nueva. Primero hay que establecer qué observación es un hecho y después formular al menos dos explicaciones en competencia; a continuación, se eligen los puntos de presupuesto, las subtareas o las ablations de arquitectura que permitan distinguir entre esas explicaciones. Si el resultado solo se sostiene en un punto, no debe afirmarse que cambia la scaling trajectory.

\begin{importantbox}{Plantilla de siete pasos: del aula al experimento}
\begin{enumerate}
  \item Muestrear y revisar a mano datos reales de entrenamiento o de evaluación;
  \item Definir una metric continua y una tarea discreta del usuario, para no quedarse con una sola puntuación;
  \item Enumerar las tareas latentes, las estrategias y las hipótesis estructurales;
  \item Ejecutar la baseline y la intervention en varios puntos de presupuesto;
  \item Hacer ablations de una sola variable sobre las estructuras clave;
  \item Reportar los intervalos de ajuste, los errores, los casos de falla y la distribución de los datos;
  \item Indicar con claridad si la conclusión es una observación, una hipótesis de mecanismo, un juicio personal o una extrapolación.
\end{enumerate}
\end{importantbox}

\subsection{Preguntas de autoevaluación}

Estas preguntas sirven para comprobar si se domina de verdad el razonamiento de las conferencias y no solo se recuerdan las conclusiones. Al responder, conviene citar fórmulas, variables de las figuras o experimentos contrafácticos.

\begin{enumerate}
  \item ¿Por qué el next-token objective puede verse como massive multi-task learning? ¿Y por qué no garantiza que todas las tareas se aprendan de forma equilibrada?
  \item Cuando la loss global desciende de manera suave, ¿por qué la exact-match accuracy puede subir de golpe?
  \item ¿Cómo se explica el U-shaped scaling mediante repeat, quote repair e instruction following?
  \item ¿Por qué no basta con que un método de investigación lleve la delantera en un solo punto de presupuesto para afirmar que es más scalable?
  \item ¿En qué condiciones el inductive bias ayuda y en cuáles se convierte en un cuello de botella?
  \item ¿Qué hipótesis elimina cada uno de los cuatro pasos de la transformación de encoder-decoder a decoder-only?
  \item ¿Por qué los datos long-input/short-target podrían ser especialmente adecuados para encoder-decoder?
  \item ¿Por qué la bidirectional attention aumenta el cómputo repetido en los diálogos de varios turnos?
  \item ¿Qué supuesto estructural impone la MLE a la generación abierta? ¿Y qué problema sigue sin resolver el RLHF?
  \item Si en el futuro la compute trend cambia por restricciones energéticas, ¿cómo debería actualizarse el método de análisis de Hyung Won?
\end{enumerate}

\subsection{Lecturas adicionales}

Esta sección enumera solo materiales originales directamente relacionados con la argumentación de la clase. Al leerlos, se recomienda reproducir los ejes y los supuestos de las figuras en lugar de limitarse a extraer las conclusiones.

\begin{itemize}
  \item Kaplan et al., \emph{Scaling Laws for Neural Language Models}: relación empírica de ley de potencias de la loss global.
  \item Wei et al., \emph{Emergent Abilities of Large Language Models}: curvas a nivel de tarea y emergent behavior.
  \item Schaeffer et al., \emph{Are Emergent Abilities of Large Language Models a Mirage?}: efecto de la metric choice en la forma de las curvas.
  \item Sutton, \emph{The Bitter Lesson}: relación a largo plazo entre los métodos generales, el cómputo y la estructura diseñada a mano.
  \item Vaswani et al., \emph{Attention Is All You Need}: el Transformer encoder-decoder original.
  \item Chung et al., \emph{Scaling Instruction-Finetuned Language Models}: evidencia de instruction tuning de la serie FLAN.
\end{itemize}

\end{document}
